\section{Introduction}

Semantic-ID (\sid) tokenization has become a practical design point for
generative recommendation. Systems such as TIGER map items into discrete code
sequences so that a sequence model can generate candidate identifiers rather
than score every catalog item directly~\cite{rajput2023tiger}. This shift has
made the tokenizer itself a first-order artifact. Recent work has expanded the
space quickly: residual quantization frameworks, recommender-native encoding,
collaborative tokenization, differentiable tokenization, non-uniform capacity,
qualified collisions, variable-length interfaces, and drift-aware refreshes all
change the item address space exposed to a generator~\cite{ju2025grid,liang2026resid,zhu2024cost,fu2026diger,wei2026card,hu2026quasid,cheng2026capsid,feng2026dact}.

The evaluation practice has not caught up with this artifact boundary. Prior
recommender tooling has argued for behavioral tests beyond NDCG-style
aggregates~\cite{chia2022reclist}, but SID tokenizers introduce a distinct
artifact: the generated item address space. A
tokenizer can support a usable downstream Recall@K while still hiding failure
modes that are hard to diagnose from aggregate ranking metrics: underused
codebooks, repeated full codes, semantically neat but behaviorally poor
prefixes, tail items sharing too little capacity, or prefix structures that
increase serving cost. These issues matter because \sid{}s are not merely
features. They define the address space exposed to the generator.

This gap is especially awkward for resource reuse. A new tokenizer repository
may expose checkpoints, item mappings, codebooks, or only intermediate feature
files. Downstream users then need to answer basic engineering questions before
training a generator: do all catalog items have valid codes, are many items
aliased to the same full code, are collisions concentrated in the tail, and do
prefix neighborhoods reflect user co-occurrence rather than only metadata
taxonomy? Today these checks are usually reimplemented ad hoc inside each
paper. The resulting evidence is difficult to compare because each method
reports a different mixture of codebook usage, downstream ranking, and
qualitative examples.

We present \tool, a resource for auditing item-to-\sid{} tokenizer artifacts
before or alongside downstream recommendation evaluation. The goal is
deliberately narrower than proposing a new tokenizer. \tool asks a concrete
question: given a public tokenizer artifact that maps items to code sequences,
what can be said about its capacity use, collisions, behavioral alignment,
head-tail allocation, and deployment cost without retraining a full generator?

The resource contribution is threefold. First, \tool defines a small
mapping-first interface for \sid assignments, item metadata, interaction
histories, and optional generator outputs. Second, it implements D1--D5a
diagnostics for utilization, collisions, semantic-collaborative alignment,
head-tail capacity, and structural serving-cost proxies, with D6 churn support
for continual-tokenizer settings. Third, it demonstrates the audit on public
GRID/RQ-KMeans-style, ReSID/GAOQ, sanity, and controlled-stressor artifacts.
The same-item Musical case study and stressor table support diagnostic claims:
they separate full-code collision pressure, interaction-qualified collision
risk, collaborative-prefix recovery, tail capacity, and structural prefix
fan-out before a downstream generator is trained. The paper follows the CIKM
Resource four-page constraint, where appendices count toward the limit and
detailed logs are therefore placed in the online artifact~\cite{cikm2026resource}.
